\documentclass[a4paper,11pt]{article}
\usepackage{a4wide}
\usepackage{geometry}
  \geometry{paper=a4paper,top=2.2cm,bottom=2.9cm}
   \geometry{right=2.3cm,left=2.3cm}
% Feynman diagrammen
\usepackage{feynmp}
\unitlength = 1mm
% Afbeeldingen
\usepackage{graphicx} % afbeeldingen
\usepackage{subfigure} % pakket voor het gebruiken van sub-figuren
\usepackage{float} % laat toe [H] te gebruiken bij plaatsen figuren
% Taal
\usepackage[dutch]{babel}
%\hyphenation{}
% De formules
\usepackage{amssymb, amsmath} % wiskunde
\numberwithin{equation}{section} % (#.#) i.p.v. (#)
% xetex
\usepackage{fontspec}
\usepackage{xunicode}
\usepackage{xltxtra}
\setromanfont[Mapping=tex-text]{Calibri}
% appendix
\usepackage{appendix}


\title{Dwarse energiestroom in de LHeC met saturatie}
\author{Ruben Van Boxem}

\begin{document}

% titelpagina
\fontsize{12pt}{14pt}\selectfont

\begin{center}

\includegraphics[height=3cm]{Afbeeldingen/UA.eps}

\vspace{1cm}

\fontsize{14pt}{17pt}\selectfont
% De Faculteit:
\textsc{Faculteit Wetenschappen} \\
\textsc{Departement Fysica}
\fontsize{12pt}{14pt}\selectfont
\vspace{0.3cm}

\vspace{1.2cm}

%Het academiejaar: aanpassen!
Academiejaar 2009--2010

\vspace{2.8cm}

\fontsize{17.28pt}{21pt}\selectfont

% De titel van de thesis:
\textsc{Dwarse energiestroom in de LHeC met saturatie}

\fontseries{m}
\fontsize{12pt}{14pt}\selectfont

\vspace{3cm}

% De auteur van de thesis:
Ruben \textsc{Van Boxem}	


\vspace{2cm}

Promotor: Prof. Dr Pierre \textsc{Van Mechelen}\\
Copromotor: Krzysztof \textsc{Kutak} \\
\vspace{2cm}
\end{center}
% De functie van de thesis:
Proefschrift voorgedragen tot \\
het behalen van de graad van\\
\textsc{Bachelor in de Fysica}


\thispagestyle{empty}
\newpage

\section*{Dankwoord}
TODO
\thispagestyle{empty}
\newpage
\fontsize{11pt}{14pt}\selectfont

\section*{Summary}
TODO
\thispagestyle{empty}
\newpage

\tableofcontents
\thispagestyle{empty}
\newpage

\section{Diepe inelastische verstrooiing}
      \paragraph{}
Om de diepere structuur van een object zoals het proton te bestuderen, kan men elektron verstrooiing gebruiken.
Door middel van een elektron met hoge energie op een proton te laten botsen kan men de onderliggende samenstelling van het proton bestuderen.
Een proces dat zich hier goed voor leent, heet DIS, oftewel \textit{Deep Inelastic Scattering}. In figuur \ref{fig:DIS} wordt een DIS proces schematisch weergegeven.
\begin{figure} [H]
  \begin{center}
    \includegraphics[width=.5\textwidth]{Afbeeldingen/placeholder.eps}
    \caption{Elektron-proton verstrooiing. Het elektron zendt een foton uit dat het proton “onderzoekt".
$W$ is de invariante massa van het uitgaande systeem (dus zonder elektron), $q$ is de overgedragen vierimpuls, $x$ (Bjorken x) is de fractie van de totale impuls van het proton gedragen door de beschouwde quark.}
   \label{fig:DIS}
  \end{center}
\end{figure}
 Men spreekt van DIS processen als $Q^2 \gg M^2$ (``diep") en $W^2 = (P+q)^2 \gg M^2$ (``inelastisch").
Omdat de waarde van $q^2$ negatief is, wordt vaak volgende variabele ingevoerd: $Q^2 :=-q^2 > 0$. De golflengte van het het foton is van de grootteorde van $1/Q$.
Dit wil concreet zeggen dat als de overgedragen impuls verhoogt, de golflengte van het foton en dus de grootte van de deeltjes waarmee het interageert, kleiner wordt.

      \paragraph{}
Men kan twee soorten DIS interacties onderscheiden: neutrale en geladen stroom (respectievelijk NC en CC). De Feynman diagrammen worden weergegeven in figuur %ZIE HIERONDER
\begin{figure} [H]
  \centering
  \subfigure[NC interactie met uitwisseling van een foton of Z-boson.]{
  \includegraphics[width=.45\textwidth]{Afbeeldingen/placeholder.eps}
  \label{fig:NC}
  }
  \subfigure[Een voorbeeld van een correct model van een kwart van een cirkelvormig vlies.]{
  \includegraphics[width=.45\textwidth]{Afbeeldingen/placeholder.eps}
  \label{fig:CC}
  }
  \caption{Soorten DIS interacties}
  \label{fig:NC-CC}
\end{figure}

      \paragraph{}
In wat volgt wordt beschreven wat er gebeurt als men $Q^2$ systematisch verhoogt.

  \subsection{Breking van schaalinvariantie}
    \paragraph{Elektron-kern verstrooiing}
Stel dat een elektron verstrooid wordt op een atoomkern (bestaande uit meerdere nucleonen).
Als $\lambda \gg d_{kern}$, ziet het foton de kern als een puntdeeltje en wordt het elastisch verstrooid op een puntdeeltje met de lading van de kern.
Als $\lambda \ll d_{kern}$, dringt het foton diep door in de kern en kan het verstrooid worden op een proton in de kern.
Men spreekt hier van diep inelastische elektron-kern verstrooiing, en men kan het volgende schrijven:
\begin{align}
x_N = \frac{M}{M_N} \left( \frac{Q^2}{2M\nu} \right) = \frac{1}{A}
\end{align}
waarbij $M_N$ de massa van de kern is, $M$ de massa van het proton, $\nu$ de frequentie van het foton en $A$ het aantal nucleonen in de kern.
Als men de werkzame doorsnede van deze processen zou uitzetten t.o.v. $x_N$, ziet men een duidelijke verschil tussen de gevallen waar $\lambda \gg d_{kern}$ en die waar $\lambda \ll d_{kern}$.
Dit fenomeen heet \textit{schaalinvariantie}, waarbij de waargenomen werkzame doorsnede blijkbaar afhangt van de grootte van de schaal, $Q^2$.

      \paragraph{Elektron-proton verstrooiing}
Zolang dat $\lambda \gg d_{proton}$, ziet het foton het proton als puntdeeltje. Het wordt elastisch verstrooid door het proton.
Als $Q^2$ zo groot is dat $\lambda \ll d_{proton}$, wordt er opnieuw ``ingezoomd", waarbij het foton nu de bouwstenen van het proton kan “zien" (i.e. ermee kan interageren).
Het foton ziet dus de 3 quarks als puntdeeltjes en men spreekt van diep inelastische elektron-proton verstrooiing.
Kinematisch geldt het volgende (zie appendix \ref{app:DIS}):
\begin{align} \label{eq:Bjorkenx}
x = \left( \frac{Q^2}{2p\cdot q} \right)
\end{align}
$x$ wordt aangeduid als de “Bjorken schaalvariabele".
De waarschijnlijkheid van zo'n proces worden dus enkel en alleen bepaalt door $x$ en niet $Q^2$ en $p\cdot q$ afzonderlijk.

      \paragraph{Breking van schaalinvariantie}
In volgorde van stijgende $Q^2$ heeft men dus voor de werkzame doorsnede:
\begin{itemize}
  \item\textit{kern}schaling met een piek op $x_N=1$
  \item  breking van schaalinvariantie
  \item \textit{proton}schaling met een piek \textit{rond} $x \sim 1$
  \item breking van schaalinvariantie
  \item \textit{Bjorken}schaling met een \textit{uitgesmeerde} piek rond $x \sim 1/3$ (3 quarks)
\end{itemize}
Als men $Q^2$ blijft vergroten, verwacht men logischerwijs nogmaals een breking van schaalinvariantie te zien, en deze worden ook geobserveerd.
De schijnbare breking van schaalinvariantie ligt in lijn met de voorspellingen van QCD, de theorie van de sterke wisselwerking.

  \subsection{Zeequarks, gluonen en $\alpha_S$}
Quarks blijken niet uit een diepere structuur te bestaan (zodat er niet opnieuw een breking van schaalinvariantie optreedt).
In plaats van een diepere quark structuur, beschrijft QCD het proton als 3 valentiequarks (de klassieke $u u d$ combinatie) plus een willekeurig aantal zeequarks ($q\bar{q}$ paren).
Het foton ziet naast drie puntdeeltjes ook een zee van quarks.
Zeequarks ontstaan uit gluonen ($g\rightarrow q\bar{q}$) die zelf uitgestraald worden door andere quarks ($q\rightarrow gq$).
Het feit dat er extra deeltjes moeten bestaan in een proton volgt ook uit de experimentele waarneming dat slechts de helft van de impuls van het proton gedragen wordt door z’n valentiequarks.
Er treedt dus breking van schaalinvariantie op omdat elk parton (quark of gluon) bij een grotere waarde van $Q^2$, een kleiner deel  $x$ van de impuls $P$ van het proton zal dragen.
Deze breking van schaalinvariantie heeft volgende vorm (zie %REFERENTIE NAAR ERGENS LATER%):
\begin{align}
\alpha_S R \ln{\left( \frac{Q^2}{\mu^2}\right)}
\end{align}
Hier is $R$ een observabele en $\mu$ de renormalisatie schaal.
Merk op dat in tegenstelling tot de vorige brekingen van schaalinvariantie, deze logaritmisch is, wat bijzonder moeilijk te meten was toen deze ontdekking werd gedaan.

  \subsection{Structuurfuncties}
De meetbare grootheid die het geheel verbindt met harde waarnemingen, zijn de structuurfuncties.
Deze waarschijnlijkheidsverdelingen in functie van $x$ en $Q^2$  geven een indicatie van de kans om in een DIS proces met een parton te botsen met die parameters.
Met de juiste definities kan men de werkzame doorsnede voor een NC proces als volgt schrijven (als $M^2/Q^2 \rightarrow 0$):
\begin{align} \label{eq:SF}
\frac{d\sigma}{dxdQ^2} &= \frac{2\pi \alpha^2}{xQ^4} \left(Y_+ F_2 \pm Y_- x F_3 - y^2 F_L \right)
\end{align}
Een uitgebreide discussie van bovenstaande uitdrukking staat in appendix \ref{app:SF}.
Als er enkel een foton kan worden uitgewisseld ($\sigma_\gamma \gg \sigma_Z$, dus de interactie via Z wordt verwaarloosd), kan men aantonen dat $F_3 \approx 0$ (zie appendix \ref{app:SF}).
Zoals in REFERENTIE aangetoond kunnen de structuurfuncties ontwikkeld worden als een reeks in $\alpha_S$.

  \subsection{$\alpha_S$}
    \subsubsection{QED en $\alpha$}
Om de vorm en het gedrag van de sterke koppeling te begrijpen, is het handig om eerst kort de koppeling in QED ($\alpha = e^2/4\pi$) te bekijken.
De processen die bijdragen aan een interactie van een foton met een enkel elektron worden in figuur \ref{fig:QED} weergegeven.
\begin{figure} [H]
  \begin{center}
    \includegraphics[width=.5\textwidth]{Afbeeldingen/placeholder.eps}
    \caption{Feynman diagrammen van de processen die bijdragen tot de interactie tussen een foton en een elektron.}
   \label{fig:QED}
  \end{center}
\end{figure}
De lading van een kaal elektron wordt afgeschermd door de “polarisatie van het vacuüm”, waardoor het foton extra $e^+ e^-$-paren “voelt”.
Hoe korter de golflengte van het foton, hoe meer $e^+ e^-$-paren het zal voelen.
Dit zorgt ervoor dat $\alpha$ toeneemt met $Q^2$.
De normalisatie van $\alpha$ komt uit vergelijking met experimentele data.
Algemeen kan men de elektromagnetische koppelingsfactor schrijven als volgt.
\begin{align}
\alpha(Q^2) = \alpha_0 \left[ 1+ \frac{\alpha_0}{3\pi} \log{\frac{Q^2}{M^2}} + \left( \frac{\alpha_0}{3\pi} \log{\frac{Q^2}{M^2}} \right)^2 +\hdots \right]
\end{align}
De \textit{loops} in de hogere orde Feynman diagrammen kunnen een oneindige bijdrage hebben indien er geen \textit{cut-off} waarde voor de


  \subsection{Dwarse energiestroom}
    \subsubsection{Formulering}
    \subsubsection{Gluon dichtheid}
    \subsubsection{$F_2$}
    \subsubsection{$\mathcal{F}_2$}
    \subsubsection{Saturatie}

\section{Resultaten}

\section{Bespreking}

\section{Besluit}

\newpage
\appendix
\section{Diep Inelastische Verstrooiing} \label{app:DIS}
Als men figuur \ref{fig:DIS} beschouwt in het IMF ($m_e \approx m_p \approx 0$), geldt wegens behoud van energie:
\begin{align}
&(xP+q)^2 = m_e = 0 \notag \\
\Leftrightarrow &x^2 P^2 + 2xPq + q^2 = 0 \notag \\
\Leftrightarrow &2xPq -Q^2=0
\end{align}
Waarbij gebruikt is dat $-q^2 = Q^2 \gg x^2P^2$ in het  IMF. Dit laatste is equivalent met \eqref{eq:Bjorkenx}.

\section{DIS werkzame doorsnede} \label{app:SF}
Uit kwantum veldentheorie kan men de werkzame doorsnede voor een NC proces zoals in figuur \ref{fig:NC} halen:
\begin{align}
\frac{d\sigma}{dxdQ^2} &= xs \frac{d\sigma}{dxdQ^2} = \frac{2\pi y \alpha^2}{Q^4} \sum_j \eta_j L_j^{\mu \nu} W_{\mu \nu}^j
\end{align}
Waarbij de som over j loopt over
\begin{itemize}
  \item $\gamma$, met $\eta_\gamma=1$,
  \item $\gamma Z$, met $\eta_{\gamma Z} = \left( \frac{G_F M_Z^2}{2\sqrt{2}\pi \alpha} \right) \left(\frac{Q^2}{Q^2+M_Z^2} \right)$,
  \item $Z$, met $\eta_Z = \eta_{\gamma Z}^2$.
\end{itemize}
Hier stelt de tweede term de foton en Z uitwisseling voor, en de interferentie tussen de twee. $G_F$ is de Fermi koppeling en $\alpha$ de QED koppeling.
Er is naast $x$ en $Q^2$ nog een derde variabele nodig, die te maken heeft met de totale energie van het systeem.
Dit kan de Mandelstam variabele s zijn, of het relatieve energieverlies van het elektron in het massamiddelpuntsysteem y, zoals hieronder gedefinieerd.
\begin{align}
y &= \frac{p \cdot q}{p \cdot k} = \frac{\Delta E}{E} ,& s = (k+p)^2 \simeq \frac{Q^2}{xy}
\end{align}
$y$ wordt wel eens als \textit{rapidity} gebruikt, wat een Lorentz invariante parameter is voor de verstrooiingshoek $\hat{\theta}$ in het massamiddelpuntstelsel.
In dit geval geldt volgende relatie:
\begin{align}
y = \frac{1}{2} \left( 1-\cos{\hat{\theta}} \right)
\end{align}
$L^{\mu \nu}$ is de gekende tensor van het lepton knooppunt in functie van $k$ en $k'$.
$W_{\mu \nu}$ is een onbekende tensor van het hadron knooppunt.
Deze laatste kan dus enkel geparametriseerd worden aan de hand van wat wel gekend is.
De algemene uitdrukking ingeval van niet gepolariseerde DIS is als volgt:
\begin{align}
W_{\mu \nu} = \left( -g_{\mu \nu} + \frac{q_\mu q_\nu}{q^2} \right) F_1(x,Q^2) + \frac{\hat{P}_\mu \hat{P}_\nu}{p \cdot q} F_2(x,Q^2) -  i\epsilon_{\mu \nu \alpha \beta} \frac{q^\alpha p^\beta}{2p \cdot q} F_3(x,Q^2)
\end{align}
Hier zijn $F_i$ de structuurfuncties, afhankelijk van $x$ en $Q^2$.
De derde term bevat een uitwendig product (de $\vec{q} \times \vec{q}$ vorm), waardoor deze term pariteit niet behoudt.
Deze term zal dan ook $\approx 0$ indien de Z bijdrage verwaarloosbaar is, omdat voor zover gekend pariteit nog steeds behouden is voor de sterke wisselwerking.
Men bekomt \eqref{eq:SF} door $L^{\mu \nu}$ in te vullen en volgende variabelen te definiëren:
\begin{align}
Y_\pm &= 1 \pm (1-y)^2 \notag \\
F_L &= F_2 - 2x F_1 \notag
\end{align}
In DIS processen, kan men vaak uitgaan van $F_L = 0$, zodat $F_2$ voldoende is om de hele werkzame doorsnede te bepalen.

\newpage
\begin{thebibliography}{99}

\bibitem{Martin}
  Alan D. Martin,
  \emph{Proton Structure, Partons, QCD, DGLAP, and Beyond}.
  Institute for Particle Physics Phenomenology,
  University of Durham
  Durham, DH1 3LE, UK
  12 mei, 2008.

\end{thebibliography}



\end{document}